\documentclass[10pt,a4paper]{amsart}
\usepackage[margin=19mm]{geometry}
\usepackage{amsmath,amssymb,mathtools}
\usepackage[scr=rsfs,cal=euler]{mathalfa}
\usepackage{xfrac}
\usepackage[unicode,hidelinks,bookmarks=false]{hyperref}
\newcommand{\ie}{i.e.\ }
\newcommand{\cf}{cf.\ }
\newcommand{\resp}{respectively\ }
\newcommand{\Subsection}{\subsection}
\providecommand{\nobreakdash}{\nobreak\hbox{-}}
\setlength{\emergencystretch}{2em}
\multlinegap0pt
\usepackage{tikz}
\usetikzlibrary{cd}
\usepackage[shortlabels]{enumitem}
\usepackage{mathtools}
\usepackage{tikz}
\usepackage{amssymb}
\newtheorem{definition}{Definition}
\newtheorem{corollary}{Corollary}
\newtheorem{lemma}{Lemma}
\newtheorem{proposition}{Proposition}
\newcommand{\Obj}{\ensuremath{\operatorname{Obj}}}
\newcommand{\Ring}{\ensuremath{\mathbb R \hspace{-.07cm} \operatorname{ing}}}
\newcommand{\lcm}{\mathrm{lcm}}
\title{The dimension of the tangent bundle and the universality of the vertical lift}
\author{Florian Schwarz}
\date{Source: arXiv:2602.15807v1 (CC BY 4.0)}
\begin{document}
\maketitle

\noindent\emph{Source and licence.} The dimension of the tangent bundle and the universality of the vertical lift. Version arXiv:2602.15807v1 (CC BY 4.0), distributed by arXiv under the Creative Commons Attribution 4.0 International licence, http://creativecommons.org/licenses/by/4.0/, as recorded in the arXiv licence metadata for this exact version and shown on the abstract page https://arxiv.org/abs/2602.15807v1. Full licence notice: \texttt{licenses/arxiv-2602.15807-CC-BY-4.0.txt}.\par\smallskip

\noindent\textit{Editorial scope.} Counted unit: the article's proposition prop:characteristic\_as\_dim\_on\_Ring (source0.tex lines 1337--1339). The article's complete original proof of it is delivered with it (source0.tex lines 1340--1423, one \texttt{proof} environment of 84 lines). No mathematics is written, repaired or re-derived: the statement and the proof are the article's own lines, in the article's own notation, and no retained line is substituted or re-flowed.  Retained uncounted as the source-local context the proof needs: lines 418--440; lines 463--473; lines 479--505; lines 1287--1289; lines 1302--1308.  Nothing was omitted from any retained span, so the package carries no omission record.  No retained line contains a citation, so the wrapper carries no bibliography and no bibliography entry.  The retained lines contain the article's own diagram code, which the wrapper typesets with the standard tikz packages; no image file is delivered and the package declares no asset dependency.  Editorial formatting only, in addition: an AMS \texttt{amsart} preamble that restates the article's own declarations and macros used by the retained lines, namely the environments \texttt{definition}, \texttt{corollary}, \texttt{lemma}, \texttt{proposition} on separate counters, together with the standard packages those lines need.  The retained environments print in this document as Definition 1, Corollary 1, Lemma 1, Proposition 1 in order of appearance, whereas the article prints the same items with the numbering of its own full document.  The complete native source of the article as distributed in the arXiv e-print for this exact version is bundled unchanged as \texttt{sources/194828/source0.tex}.
\begin{definition}\label{def:monoid-valued-dimension}
    Let $M$ be a commutative monoid and $\mathbb X$ be a category. Then an \textbf{$M$-valued dimension} on $\mathbb X$ is a function $\dim : \pi_0(\mathbb X) \to M$ from the set of isomorphism classes of objects of $\mathbb X$ to the monoid $M$, satisfying the following property:\\
        For any pullback 
        % https://q.uiver.app/#q=WzAsNCxbMSwwLCJBIl0sWzEsMSwiQiJdLFswLDEsIkMiXSxbMCwwLCJBIFxcdGltZXNfQiBDIl0sWzMsMF0sWzAsMSwicl8yIl0sWzIsMSwicl8xIiwyXSxbMywyXV0=
    \[\begin{tikzcd}
    	{A \times_B C} & A \\
    	C & B
    	\arrow[from=1-1, to=1-2]
    	\arrow[from=1-1, to=2-1]
    	\arrow["{r}", from=1-2, to=2-2]
    	\arrow["{f}"', from=2-1, to=2-2]
    \end{tikzcd}\]
        of $f$ along a retraction $r$,  
        % in which $f$ is either a retraction or else a section, 
    \begin{itemize}
        \item if $f$ is a retraction, or
        \item if $f$ is a section,
    \end{itemize} the equation
        $$
        \dim(A \times_B C) + \dim(B) = \dim(A) + \dim(C)
        $$
        holds.
\end{definition}

\begin{corollary}\label{cor:adjunctions_full_subcats_and_dimensions}
    Let $\mathbb X$ and $\mathbb Y$ be categories, let $\mathbb X$ have an $R$-valued dimension $\dim_\mathbb X: \pi_0(\mathbb X) \to R$. 
    \begin{enumerate}[label = (\alph*)]
        \item Let $F: \mathbb Y \to \mathbb X$ be a functor which is a right adjoint. Then the assignment that sends $Y \in \pi_0(\mathbb Y)$ to
    $$
    Y \mapsto \dim_\mathbb X(F(Y))
    $$
    is an $R$-valued dimension on $\mathbb Y$.
    \item     Let $\mathbb Y$ be a full subcategory of $\mathbb X$. Let $\dim:  \pi_0(\mathbb X) \to R$ be a $R$-valued dimension on $\mathbb X$. Then $\dim|_{\Obj(\mathbb Y)}$ is an $R$-valued dimension on $\mathbb Y$.
    \end{enumerate}
\end{corollary}

\begin{definition}\label{def:tangent_dimension}
    Let $M$ be a commutative monoid and $(\mathbb X , T)$ be a tangent category. Then an \textbf{$M$-valued tangent dimension} on $(\mathbb X, T)$ is a function $\dim : \pi_0(\mathbb X) \to M$ satisfying the following property:\\
    For any $T$-pullback 
        % https://q.uiver.app/#q=WzAsNCxbMSwwLCJBIl0sWzEsMSwiQiJdLFswLDEsIkMiXSxbMCwwLCJBIFxcdGltZXNfQiBDIl0sWzMsMF0sWzAsMSwicl8yIl0sWzIsMSwicl8xIiwyXSxbMywyXV0=
    \[\begin{tikzcd}
    	{A \times_B C} & A \\
    	C & B
    	\arrow[from=1-1, to=1-2]
    	\arrow[from=1-1, to=2-1]
    	\arrow["{r}", from=1-2, to=2-2]
    	\arrow["{f}"', from=2-1, to=2-2]
    \end{tikzcd}\]
        along a retraction $r$,
        \begin{itemize}
        \item if $f$ is a retraction, or
        \item if $f$ is a section,
        \end{itemize} the equation
        $$
        \dim(A \times_B C) + \dim(B) = \dim(A) + \dim(C)
        $$
        holds.
    An $M$-valued tangent dimension is \textbf{strong} if $M$ is a commutative rig and there is an element $a \in \mathbb M$ such that for every object $X \in \mathrm{Ob}(\mathbb X)$ the equation
    $$
    \dim (T(X)) = a \cdot \dim(X)
    $$
    holds.
\end{definition}

\begin{lemma}\label{lem:maximal_order_elem}
Let $(R,+,\cdot,0)$ be a ring of characteristic $n$. Then there exists an element $x_R \in R$ such that $k \cdot x_R \neq 0$ for all $k < n$. We call this element a \textbf{maximal order element}.
\end{lemma}

\begin{lemma}\label{lem:characteristic_section_retraction}
    Given a section $A \xrightarrow{s} B $ with retraction $ B \xrightarrow{r} A$ in $\Ring_n$, the inequality
    $$
    \mathrm{char}(A) \leq \mathrm{char}(B)
    $$
    holds. Additionally, $\mathrm{char}(B)$ is a multiple of $\mathrm{char}(A)$.
\end{lemma}

\begin{proposition}\label{prop:characteristic_as_dim_on_Ring}
    The characteristic $\mathrm{char}$ is an $(\mathbb N_\infty^*, \lcm )$-valued dimension on $\Ring_n$, $\Ring_1$ and on $\Ring_u$. 
\end{proposition}

\begin{proof}
    We will first check the criteria of Definition \ref{def:monoid-valued-dimension} for $\Ring_u$ and then for $\Ring_1$.

    In $\Ring_u$, given any morphism $f: R \to R'$, the characteristic can not increase $\mathrm{char}(R')\leq \mathrm{char}(R) $ since $ 0 = f(0) = f(\mathrm{char}(R) \cdot 1)$. Thus, given a section retraction pair between $R$ and $R'$, $\mathrm{char}(R) = \mathrm{char}(R')$. Now let 
            % https://q.uiver.app/#q=WzAsNCxbMSwwLCJBIl0sWzEsMSwiQiJdLFswLDEsIkMiXSxbMCwwLCJBIFxcdGltZXNfQiBDIl0sWzMsMF0sWzAsMSwicl8yIl0sWzIsMSwicl8xIiwyXSxbMywyXV0=
    \[\begin{tikzcd}
    	{A \times_B C} & A \\
    	C & B
    	\arrow[from=1-1, to=1-2]
    	\arrow[from=1-1, to=2-1]
    	\arrow["{r}", from=1-2, to=2-2]
    	\arrow["{f}"', from=2-1, to=2-2]
    \end{tikzcd}\]
    be any of the pullbacks from Definition \ref{def:monoid-valued-dimension} in $\Ring_u$.
    Then $f$ and $r$ are both parts of section-retraction pairs and the characteristic of all 4 vertices is the same. Therefore
    $$
    \lcm(\mathrm{char}(A), \mathrm{char}(C)) = \lcm(\mathrm{char}(A \times_B C) , \mathrm{char}(B))
    $$
    holds trivially.
    


    
    
    Now let 
    % https://q.uiver.app/#q=WzAsNCxbMSwwLCJBIl0sWzEsMSwiQiJdLFswLDEsIkMiXSxbMCwwLCJBIFxcdGltZXNfQiBDIl0sWzMsMF0sWzAsMSwicl8yIl0sWzIsMSwicl8xIiwyXSxbMywyXV0=
    \[\begin{tikzcd}
    	{A \times_B C} & A \\
    	C & B
    	\arrow[from=1-1, to=1-2]
    	\arrow[from=1-1, to=2-1]
    	\arrow["{r}", from=1-2, to=2-2]
    	\arrow["{f}"', from=2-1, to=2-2]
    \end{tikzcd}\]
    be any of the pullbacks from Definition \ref{def:tangent_dimension} in $\Ring_n$. 
    Assuming that $A$ and $C$ have finite characteristic, let $x_A \in A$ and $x_C \in C$ be the maximal order elements constructed in Lemma \ref{lem:maximal_order_elem}.
    We will distinguish the cases when $f$ is a section and when $f$ is a retraction.

    \textit{Case 1:} Suppose $f$ is a retraction, in particular it is surjective. Then there is a $y_C \in C$ such that $f(y_c) = r(x_A)$ and thus we have $(x_a, y_C) \in A \times_B C$ such that for $k\leq \mathrm{char}(A)$, $k \cdot (x_A, y_C) = (k \cdot x_A , k \cdot y_C) \neq 0$. Thus $\mathrm{char}(A \times_B C)$ is a multiple of $\mathrm{char}(A)$.

    Analogously there is an $y_A \in A$ such that $(y_A, x_C) \in A \times_B C$ and thus $\mathrm{char}(A \times_B C)$ is a multiple of $\mathrm{char}(C)$. This shows that $\mathrm{char}(A \times_B C)$ is a multiple of $\lcm(\mathrm{char}(A), \mathrm{char}(C))$.
    Since $A \times_B C$ is a subset of $A \times C$, the equation
    $$
    \mathrm{char}(A \times_B C) = \lcm(\mathrm{char}(A), \mathrm{char}(C))
    $$
    holds. Due to Lemma \ref{lem:characteristic_section_retraction}, $\mathrm{char}(B)$ divides $\mathrm{char}(A)$, implying that
    $$
    \lcm(\mathrm{char}(A \times_B C), \mathrm{char(B)}) = \lcm( \lcm(\mathrm{char}(A), \mathrm{char}(C)),\mathrm{char}(B)) = \lcm(\mathrm{char}(A), \mathrm{char}(C)).
    $$
    This is exactly the equation we needed to show for Definition \ref{def:monoid-valued-dimension}.

    \textit{Case 2:} Suppose $f$ is a section. Then, due to Lemma \ref{lem:characteristic_section_retraction}, $\mathrm{char}(B)$ and $\mathrm{char}(C)$ both divide $\mathrm{char}(A)$. 
    
    Let $s: B \to A$ be a section for the retraction $r$. We can express the maximal order element $x_A \in A$ as $x_A = s(r(x_a))+ z_A$ with $r(z_A) = 0$, by choosing $z_A = x_A - s(r(x_A))$.

    For $k = \lcm(\mathrm{char}(A \times_B C), \mathrm{char}(B))$, we observe that 
    $$
    k \cdot x_A = k \cdot s(r(x_A)) + k \cdot z_A.
    $$
    Since $k$ is a multiple of $\mathrm{char}(B)$, $k \cdot s(r(x_A)) = s(k \cdot r(x_A)) = 0$. 
    Since $r(z_A) = 0$, we have the element $(z_A ,0) \in A \times_B C$ and since $k$ is a mutiple of $\mathrm{char}(A \times_B C)$, $k \cdot (z_A,0) = 0$, implying that $k \cdot z_A =0$.

    Thus $k \cdot x_A = 0$ and since $x_A$ had maximal order this shows that $\mathrm{char}(A)$ divides $\lcm(\mathrm{char}(A \times_B C), \mathrm{char}(B))$. Since $\mathrm{char}(C)$ divides $\mathrm{char}(A)$, the equation $\lcm(\mathrm{char}(A), \mathrm{char}(C)) = \mathrm{char}(A)$ holds, implying that $\lcm(\mathrm{char}(A), \mathrm{char}(C)) = \mathrm{char}(A)$ divides $\lcm(\mathrm{char}(A \times_B C), \mathrm{char}(B))$. Since $A \times_B C$ is a subset of $A \times C$, this shows that 
    $$
    \lcm(\mathrm{char}(A \times_B C), \mathrm{char}(B)) = \lcm(\mathrm{char}(A), \mathrm{char}(C)).
    $$
    This is exactly the equation we needed to show for Definition \ref{def:monoid-valued-dimension}.

    Now it only remains to show that the equation 
    $$
    \lcm(\mathrm{char}(A \times_B C), \mathrm{char}(B)) = \lcm(\mathrm{char}(A), \mathrm{char}(C)).
    $$
    also holds when $\mathrm{char}(A)$ or $\mathrm{char}(C)$ are infinite. If $\mathrm{char}(B)=\infty$, it holds trivially. Suppose $\mathrm{char}(B)$ is finite. Since in the case that $f$ is a section $\mathrm{char}(C) \leq \mathrm{char}(B)$, and in the case that $f$ is a retraction, $C$ and $A$ can be swapped, we can assume without loss of generality that $\mathrm{char}(A) = \infty$. 

    Let $n \in \mathbb N\setminus\{0\} $ be arbitrary. Let $x_n \in A$ be such that $n \cdot \mathrm{char}(B) \cdot x_n \neq 0$. Then $r(\mathrm{char}(B) \cdot x_n) = 0$ and thus $(\mathrm{char}(B) \cdot x_n, 0) \in A \times_B C$ fulfills that 
    $$
    n \cdot (\mathrm{char}(B) \cdot x_n, 0) \neq 0.
    $$
    Since $n$ was arbitrary, this shows that $\mathrm{char}(A \times_B C) = \infty$ and thus 
    $$
    \lcm(\mathrm{char}(A \times_B C), \mathrm{char}(B)) = \lcm(\mathrm{char}(A), \mathrm{char}(C)).
    $$
    holds. This concludes the proof that $\mathrm{char}$ is a $(\mathbb N^*_\infty, \lcm)$-valued dimension on $\Ring_n$. Since $\Ring_1$ is a full subcategory of $\Ring_n$, Corollary \ref{cor:adjunctions_full_subcats_and_dimensions} implies that it also is a $(\mathbb N^*_\infty, \lcm)$-valued dimension on $\Ring_1$.
\end{proof}

\end{document}
